\documentclass[12pt,a4paper]{article}
\usepackage[utf8]{inputenc}
\usepackage[spanish]{babel}

% Definimos el contador global para introducciones
\newcounter{intro}

\begin{document}
	
	% --- TÍTULO GENERAL (sin número, centrado) ---
	\begin{center}
		\textbf{ TÍTULO GENERAL DE LAS INTRODUCCIONES}
	\end{center}
	
	\section*{Título General de lasaaa Introducciones}
	\addcontentsline{toc}{section}{Título General de las Introducciones}
	
	
	% Primera introducción
	\stepcounter{intro}
	\section*{\theintro\ Introducción}
	\addcontentsline{toc}{section}{\theintro\ Introducción}
	Aquí va el texto de la primera introducción.
	
	
	\section*{Título General de las Introducciones}
	\addcontentsline{toc}{section}{Título General de las Introducciones}
	
	% Segunda introducción
	\stepcounter{intro}
	\section*{\theintro\ Introducción}
	\addcontentsline{toc}{section}{\theintro\ Introducción}
	Aquí va el texto de la segunda introducción.
	
	% Una sección normal del documento
	\section{Marco Teórico}
	Texto del marco teórico.
	
	% Tercera introducción (más adelante en el documento)
	\stepcounter{intro}
	\section*{\theintro\ Introducción}
	\addcontentsline{toc}{section}{\theintro\ Introducción}
	Aquí va el texto de la tercera introducción.
	
\end{document}